%%%%%
%
\section{Extraction de "colonnes" cons\'ecutives d'un vecteur 1D \texttt{Int\_16}}

\itemstyle{Op\'erateur~:} {\sf Extr\_tab}
\index{Extr\_tab}

%%%%%OBJET%%%%
\itemstyle{Objet~:}
 Extrait N �colonnes� cons\'ecutives d'un vecteur 1D
de type \texttt{Int\_16},  et les range dans un nouveau vecteur 1D
du m�me type.

%%%%%ARGUMENTS%%%%
\bigskip
%%%%%ENTREES%%%%
\begin{tabarg}
\multicolumn{3}{l}{\titretab {Entr\'ees}}\\
\hline
$X$ & Vecteur, {\tt Int\_16[$Dx$]} & Param\`etre 1D d'entr\'ee (unit� indiff�rente)\\
$Dx$ & Coeff., {\tt Int} & Taille de la dimension hors-temps de $X$ (nombre de �colonnes�)\\
$Dz$ & Coeff., {\tt Int} & Taille de la dimension (hors temps) du param�tre r�sultat (nombre de �colonnes�)\\
$C_{deb}$ & Coeff., {\tt Int} & Indice de d�but de l'extraction (convention "C", d�but � z�ro)\\
\end{tabarg}

\medskip
%%%%%SORTIES%%%%
\begin{tabarg}
\multicolumn{3}{l}{\titretab {Sorties}}\\
\hline
$Z$ & Vecteur, {\tt Int\_16[$Dz$]} & Sous-tableau (m\^eme unit\'e que $X$)\\
\end{tabarg}


%%%%%FORMULE%%%%
\itemstyle{Formule}

%%Extr_tab
\[
Z_k = X_{k + C_{deb}}  \qquad k=0, Dz - 1
\]


%%%% COMMENTAIRE %%%%
\itemstyle{Commentaire~:}

Lorsque l'argument $Dz$ (Le nombre de �colonnes� � extraire) 
est �gal � 1, le vecteur r�sultat est un vecteur �1D� dont
la taille de la dimension est 1, ce qui est parfaitement
�quivalent � un vecteur �scalaire�.

%%%%%AUTEUR%%%%
\itemstyle{Auteur~:}
CNRM/GMEI/TRAMM

%%%%%REFERENCE%%%%
\itemstyle{R\'ef\'erences~:}
